\documentclass[letterpaper,10pt]{article}

\usepackage[margin=0.47in]{geometry}
\usepackage[hidelinks]{hyperref}
\usepackage{titlesec}
\usepackage{enumitem}
\usepackage[english]{babel}

\IfFileExists{glyphtounicode.tex}{
  \input{glyphtounicode}
  \ifdefined\pdfgentounicode
    \pdfgentounicode=1
  \fi
}{}

\pagestyle{empty}
\raggedright
\setlength{\parindent}{0pt}
\setlength{\tabcolsep}{0pt}

\titleformat{\section}
  {\normalsize\bfseries}
  {}
  {0em}
  {}
  [\titlerule]

\titlespacing{\section}{0pt}{5pt}{3pt}

\setlist[itemize]{
  leftmargin=0.16in,
  itemsep=1pt,
  topsep=1pt,
  parsep=0pt,
  partopsep=0pt
}

\newcommand{\resumeSubheading}[4]{
  \textbf{#1} \hfill \textbf{#2} \\
  \textit{#3} \hfill \textit{#4} \\
}

\begin{document}

\begin{center}
    {\Large \textbf{Kirolos Sedra}} \\
    \textbf{RAN Data Collection Engineer / Drive Tester} \\
    Waterloo, Ontario \textbar{} Open to Montreal and relocation \textbar{}
    +1 519-580-9685 \textbar{}
    \href{mailto:kasedra@proton.me}{kasedra@proton.me} \textbar{}
    \href{https://linkedin.com/in/kirolossedra}{linkedin.com/in/kirolossedra}
\end{center}

\vspace{-6pt}

\section{Professional Summary}

RAN data collection and wireless systems engineer with hands-on experience executing \textbf{LTE/5G} field and lab measurement campaigns, including mobility testing, timestamped RF KPI collection, GPS-aligned walk tests, log validation, and post-processing. Built custom cellular data-collection workflows using \textbf{Quectel AT/QMI}, \textbf{srsRAN/Open5GS logs}, Python/Linux automation, USRP B210, and RF/network instrumentation.

\section{Experience}

\resumeSubheading
{Wireless Sensors and Devices Lab, University of Waterloo}{Sep 2024 -- Aug 2026}
{Research Assistant}{Full-time, On-site}
\begin{itemize}
    \item Built and validated a private \textbf{5G SA testbed} using an \textbf{srsRAN gNB}, \textbf{Open5GS 5GC}, \textbf{USRP B210}, and \textbf{Quectel RM502Q-AE UE}; resolved AMF/SMF/UPF registration, PDU-session, and user-plane connectivity issues.
    \item Executed attenuation-based SNR ladders, static OTA tests, mobility scenarios, and \textbf{4G/5G walk-test data collection}; measured and logged \textbf{RSRP, RSRQ, SINR, RSSI, CQI, MCS, UL/DL throughput}, and latency under load, and aligned timestamped radio/performance logs with \textbf{GPS} data for street-level analysis.
    \item Built a custom cellular/RF data-collection workflow using \textbf{Quectel AT/QMI interfaces}, srsRAN gNB and Open5GS logs, \textbf{Python/Linux automation}, tcpdump/Wireshark, and post-processing/visualization to validate data quality and investigate RF, mobility, and network issues.
    \item Characterized indoor and outdoor \textbf{5G/LTE and Wi-Fi performance for mobile-robot teleoperation}, correlating radio-quality measurements with P99 latency and throughput to quantify mobility, interference, and LTE-anchoring effects; documented results for technical and industry stakeholders.
    \item Automated wireless validation for up to \textbf{30 DUTs}, reducing manual intervention by approximately \textbf{70\%} through RF logging, network-service automation, SSH-based configuration, and PHY/transport analysis.
\end{itemize}

\resumeSubheading
{Aquasensing}{Sep 2024 -- Present}
{Firmware Specialist / Research Intern}{Part-time}
\begin{itemize}
    \item Validated and debugged production \textbf{BLE--LTE IoT gateways} through RF-informed and anechoic-chamber testing, component stubs, and cross-layer isolation of firmware, connectivity, and cloud failures.
    \item Managed cellular IoT connectivity through \textbf{Hologram MVNO} and \textbf{Bell Jasper}; debugged \textbf{nRF/Zephyr} firmware and automated provisioning and validation workflows for \textbf{70 IoT gateways}.
\end{itemize}

\resumeSubheading
{University of Waterloo}{Feb 2025 -- Apr 2026}
{Teaching Assistant, ECE 610 Broadband Networks and ECE 358 Computer Networks}{Part-time}

\resumeSubheading
{Ejada Systems}{Sep 2023 -- Aug 2024}
{Software Engineer -- Java Spring Boot, REST APIs, SQL Optimization}{Full-time}

\section{Publications}

\textbf{CCECE 2026:} Trajectory-Based Performance Characterization of Wi-Fi for Outdoor Mobile Robot Teleoperation; Scalability Analysis of Wi-Fi Broadcast Mechanisms Using Simulation and Hybrid Approaches.

\section{Education}

\textbf{University of Waterloo} -- Master of Applied Science, Electrical and Computer Engineering \hfill Sep 2024 -- Aug 2026 \\
Grade: 92\% (A+) \\
\textbf{Ain Shams University} -- Bachelor of Science, Computer Engineering \hfill Sep 2018 -- Jul 2023 \\
CGPA: 3.12/4.00

\section{Technical Skills}

\textbf{Cellular/RAN:} 5G NR SA/NSA, LTE, RAN architecture, PHY/MAC/RRC procedures, UE registration, PDU sessions, UL/DL analysis, mobility/field testing, 3GPP TS 38-series exposure \\
\textbf{Data Collection/Diagnostics:} Quectel AT/QMI, srsRAN gNB/Open5GS logs, GPS-aligned logging, Python/Linux automation, tcpdump, Wireshark, iw, nmcli, RF data post-processing and visualization \\
\textbf{RF/Performance:} RSRP, RSRQ, SINR, RSSI, SNR, CQI, MCS, throughput, latency under load, attenuators, OTA, spectrum analyzer, TinySA, VNA, anechoic chambers, interference/coexistence \\
\textbf{Simulation/Systems:} ns-3, CTTC 5G-LENA, NR-U, GNU Radio, MATLAB, Simulink, Python, C/C++, Linux, TCP/IP, SSH

\end{document}
